% Part: normal-modal-logic
% Chapter: filtrations
% Section: S5-decidable

\documentclass[../../../include/open-logic-section]{subfiles}

\begin{document}

\olfileid[pt]{nml}{fil}{dec}

\olsection{\Log{S5} é decidível}
A propriedade do modelo finito nos dá uma maneira simples de mostrar que
sistemas de lógica modal dados por esquemas são \emph{decidíveis} (O.s.\@,
que há um procedimento computável para determinar se !!a{formula} é ou não
!!{derivable} no sistema).

\begin{thm}
  \Log{S5} é decidível.
\end{thm}

\begin{proof}
  Dada $!A$, suponha que as !!{propositional variable}s que ocorrem em
  $!A$ estejam entre $p_1$, \dots, $p_k$. Como, para cada $n$, há apenas
  um número finito de modelos com $n$ mundos que atribuem valores a
  $p_1$, \dots, $p_k$, podemos enumerar, \emph{em paralelo}, todos os
  teoremas de \Log{S5}, gerando demonstrações de maneira sistemática, e
  todos os modelos que contêm $1$, $2$, \dots mundos, verificando se $!A$
  falha em algum mundo de um desses modelos. Por fim, um dos dois processos
  paralelos dará uma resposta, pois, por \olref[com][fra]{thm:generaldet}
  e \olref[fmp]{cor:S5fmp}, ou $!A$ é !!{derivable}, ou falha em um modelo
  universal finito.
\end{proof}

A demonstração acima funciona para \Log{S5} porque filtrações de modelos
universais são automaticamente universais. O mesmo vale para reflexividade
e serialidade, mas outras propriedades exigem mais trabalho.

\end{document}
